% LSPV-1 REPORT::ARC-SOURCE-INJECTIVITY-FIREWALL163::2026-10-25T01:55:00-07:00
\documentclass[11pt]{article}
\usepackage[margin=.9in]{geometry}
\usepackage{amsmath,amssymb,amsthm,xurl,hyperref}
\newtheorem{theorem}{Theorem}
\newtheorem{corollary}[theorem]{Corollary}
\title{ARC-SOURCE-INJECTIVITY-FIREWALL163\\
\large Exact Lossless-Child Counting and the Necessity of a Semantic Quotient}
\author{LSPV-1 ARC/FLMV lead}
\date{25 October 2026, America/Los\_Angeles}
\begin{document}\maketitle
\begin{center}
\textbf{PROVED/EXACT elementary encoding lemma, peer review pending.}\\
\textbf{CONDITIONAL only as an ARC application. ARC-EDGE-90 and $P$ versus $NP$ OPEN.}
\end{center}

\section{Exact target, conjectural routes, and quantifiers}
The target is the missing ARC-EDGE-90 edge for a single same-machine
uniform total deterministic exact ordinary-SAT decider, with
simultaneous charged time/space resources in the FLMV indexed,
read-only-input random-access TM model.
Use the fixed near-linear-length tagged FLMV-hard reduction
$\widehat R$ of \texttt{ARC-COOK-TAG-RECOVERY162}:
\[
 \widehat R:\{0,1\}^n\to\{0,1\}^{N(n)},\qquad
 N(n)=n\log^{O(1)}n=n^{1+o(1)}.
\]
The source tag makes $\widehat R$ injective and polylogarithmically
source-readable on the image promise.

Proposals recorded before consultation with the pre-existing
source-identity literature:
H1: a standalone child which recovers all $n$ source bits must be
near-parent-sized. PROVED for worst-case child length.
H2: a *bundle* of short children may jointly retain the entire source
while no individual child does. PROVED by elementary partition example.
H3: thus any proposed $N^{a}$-by-$N^\lambda$ source-lossless
bundle with $a+\lambda<1$ and sublinear metadata is impossible.
PROVED as a representation claim.
H4: this forces all SAT solvers to use a power-long child or
establishes $P\ne NP$. FALSIFIED as an inference; a one-bit SAT
answer is not lossless source recovery.

\section{A strict full-source-retention firewall}
\begin{theorem}[No polynomial-exponent shrink with lossless standalone children]
Let for each $n$ a deterministic encoder $C_n$ map all
$x\in\{0,1\}^n$ to a binary string $y_x$ of maximum length $m(n)$.
Suppose a fixed decoder (which may know $n$) satisfies
$D(n,y_x)=x$ for EVERY source $x$. Then
$m(n)\ge n-1$.
In particular for a near-linear-length parent
$N(n)=n^{1+o(1)}$, no such encoder can have
$m(n)\le N(n)^{\lambda+o(1)}$ for any fixed $\lambda<1$.
If output strings are fixed length $m(n)$, the stronger $m(n)\ge n$
holds.
\end{theorem}
\begin{proof}
The decoder condition makes $x\mapsto y_x$ injective.
All strings of length at most $m$ number
$\sum_{j=0}^{m}2^j=2^{m+1}-1$.
Thus $2^n\le 2^{m(n)+1}-1$ and
$m(n)\ge n$ for integer $m,n\ge0$ in this literal
unrestricted variable-length count, in particular
$m(n)\ge n-1$. The fixed-length bound is immediate.
Since $N(n)=n^{1+o(1)}$, every fixed
$\lambda<1$ gives $N(n)^{\lambda+o(1)}=o(n)$,
a contradiction to the preceding cardinality bound.
\end{proof}
The weaker $n-1$ statement deliberately absorbs any harmless
fixed delimiter convention; the raw cardinality calculation
with no side bits actually gives $m(n)\ge n$ for integer lengths.

\paragraph{Critical hypothesis.}
If the alleged short child contains only a source
\emph{pointer} and decoder $D$ may also query an
external, uncharged $n$-bit parent input, $x\mapsto y_x$
need not be injective. That violates the standalone
readable-child premise. In an ARC simulation all parent
probes, index calculations, physical or virtual input
presentation and retained ancestor memory must be paid.
This theorem never excludes fully charged parent-backed
virtual children.

\section{Lossless bundles and their exact fanout budget}
\begin{theorem}[Multi-child total-description bound]
Suppose a deterministic encoding of every $n$-bit $x$ emits
a self-delimiting binary transcript consisting of auxiliary
metadata of at most $b(n)$ bits and at most $q(n)$ complete
child strings, each of at most $m(n)$ bits, and a fixed
decoder recovers $x$ from the transcript alone. Then the
\emph{total maximum transcript length} must be at least
$n$ bits. Under any canonical child-list format using
$O(q\log(m+2)+\log(q+2))$ framing bits, necessarily
\[
 n\ \le\
 b(n)+q(n)\bigl(m(n)+O(\log(m(n)+2))\bigr)
 +O(\log(q(n)+2)).
\tag{1}
\]
Consequently if $N(n)=n^{1+o(1)}$, all metadata has
length $N^{\rho+o(1)}$ with $\rho<1$, children satisfy
$m(n)\le N^{\lambda+o(1)}$, and number
$q(n)\le N^{a+o(1)}$ with $a+\lambda<1$, the entire
lossless-bundle premise is impossible.
\end{theorem}
\begin{proof}
Since decoding recovers every x, there are at least $2^n$
distinct serialized transcripts. If the maximum
total transcript size is $M$, at most $2^{M+1}-1$
binary strings of length at most $M$ exist, so $M\ge n$.
The declared framing convention supplies the RHS of (1)
as an upper bound on $M$. The exponent conditions make
that RHS $o(n)$, contradiction.
\end{proof}
An explicit positive control partitions x into
$q=\lceil n^{1-\lambda}\rceil$ disjoint chunks of length
at most $\lceil n^\lambda\rceil$ and puts each in a
different child. The bundle recovers x and saturates
the $a+\lambda\ge1$ accounting threshold.
This is a purely information-preserving decomposition;
it supplies no satisfiability semantics. In the
illustrative ARC ledger $a=3/5$ and $\lambda=1/2$,
we have $a+\lambda=11/10>1$.
Therefore the bundle-counting theorem does NOT block
that proposed ARC parameter region.

\section{What this contributes and what remains open}
The tagged hard-image source can be cheaply recovered as
a whole if one is prepared to read all its n source
bits, but retaining that whole source independently in
each purported power-short \emph{standalone}
ordinary-SAT child cannot work. A legitimate candidate
must permit source lossiness in each individual child,
prove exact both-label parent semantics and original-name
compatibility, or transmit fragments across many
children and pay total fanout and gluing.
The source labels may be computed by a one-bit
algorithm that never reconstructs x; the counting
arguments are about declared injective \emph{encodings},
not Turing-machine lower bounds for SAT.
Nothing here implies the tagged hard image is
self-reducible at favorable cost.

\paragraph{CORPUS AND DUPLICATION AUDIT.}
Reused the freshly verified LSPV-1
\texttt{ARC-COOK-TAG-RECOVERY162} tagged construction
and the earlier image endpoint
\texttt{ARC-SINGLE-HARD-IMAGE161}. Reviewed root
\texttt{main.tex} ARC DAG, \texttt{main13.tex},
the five roles' current scoped/trace/guest-slope
reports, \texttt{MODELBRIDGE-0011} downward
oracle-elimination and \texttt{CONTRACTION-ACQUISITION1}.
The counting principle is elementary and does not
purport to be a new primary literature theorem or
generalize previous lossiness/collision entropy
beyond its stated injective-source premise.
No fresh complete repository census was performed.

\paragraph{WHAT CHANGED RELATIVE TO ROOT MAIN.}
A source-preserving \emph{standalone-child}
strategy is now explicitly ruled out at fixed
fractional exponent. This is more restricted
than the general missing ARC contractor and
does NOT upgrade any ordinary-SAT exponent pair.

\paragraph{FIVE ROLE HANDOFF.}
Role 2: design truly lossier A-dependent, semantically valid
ordinary children without off-child full-source replay or
uncharged parent calls.
Role 3: audit which original-source information is
mathematically sufficient for both SAT and UNSAT
gluing over child pairs and grandchildren.
Role 4: measure transmitted bytes, ancestor/parent
backing reads, per-probe virtual regeneration,
peak space and changed/unchanged writes.
Role 5: treat any source-identity argument which
assumes injectivity of a lossily projected child
as suspect; independently audit the framing count.
Keep exactly the existing five :00/:12/:24/:36/:48
enabled agents and all peer ownership/schedules intact.

\paragraph{Status.}
PROVED/EXACT elementary counting, peer check pending;
CONDITIONAL ARC application; OPEN actual
A-dependent typed ARC-EDGE-90, and $P$ versus $NP$.
\begin{thebibliography}{3}
\bibitem{flmv} L. Fortnow, R. Lipton,
D. van Melkebeek and A. Viglas,
\emph{Time-Space Lower Bounds for Satisfiability},
J. ACM 52(6) (2005), 835--865,
\url{https://lance.fortnow.com/papers/files/tst.pdf}.
\bibitem{tag} LSPV-1,
\emph{ARC-COOK-TAG-RECOVERY162}, 25 October 2026.
\bibitem{model} \emph{MODELBRIDGE-0011},
\path{research/review/outbox/model_bridge/2026-10-11_1300_MODELBRIDGE-0011.md}.
\end{thebibliography}
\end{document}
% END_OF_VERIFIED_REPORT::LSPV-1::ARC-SOURCE-INJECTIVITY-FIREWALL163
